\chapter{2005年北京卷（春文）}
\section{单选题}
\begin{problem}
\(\i-2\) 的共轭复数是（\(\quad\)）

\choices
    {\(2+\i\)}
    {\(2-\i\)}
    {\(-2+\i\)}
    {\(-2-\i\)}
\end{problem}

\begin{answer}
D
\end{answer}

\begin{solution}
因为 \(\i-2=-2+\i\)，所以其共轭复数为 \(-2-\i\)，故选 D.
\end{solution}

\begin{problem}
函数 \(y=|\log_2x|\) 的图象是（\(\quad\)）

\choices
    {\choicebitmap[width=0.15\paperwidth]{img/2005/beijing_liberal_spring/q2_optA.png}}
    {\choicebitmap[width=0.15\paperwidth]{img/2005/beijing_liberal_spring/q2_optB.png}}
    {\choicebitmap[width=0.15\paperwidth]{img/2005/beijing_liberal_spring/q2_optC.png}}
    {\choicebitmap[width=0.15\paperwidth]{img/2005/beijing_liberal_spring/q2_optD.png}}
\end{problem}

\begin{answer}
A
\end{answer}

\begin{solution}
函数的定义域为 \(x>0\)，且
\[
y=|\log_2x|=\begin{cases}
-\log_2x,&0<x<1,\\
\log_2x,&x\geq 1.
\end{cases}
\]
因此图象在点 \((1,0)\) 处取得最小值，左支在 \(0<x<1\) 上单调递减，右支在 \(x>1\) 上单调递增，故选 A.
\end{solution}

\begin{problem}
下列命题中，正确的是（\(\quad\)）

\choices
    {经过不同的三点有且只有一个平面}
    {分别在两个平面内的两条直线一定是异面直线}
    {垂直于同一个平面的两条直线是平行直线}
    {垂直于同一个平面的两个平面平行}
\end{problem}

\begin{answer}
C
\end{answer}

\begin{solution}
三点共线时，经过这三点的平面不唯一，所以命题 A 错误；分别位于两个平面内的两条直线可能相交、平行或异面，所以命题 B 错误；垂直于同一个平面的两条直线都与该平面的法线平行，因而两直线平行，所以命题 C 正确；垂直于同一个平面的两个平面可能相交，所以命题 D 错误，故选 C.
\end{solution}

\begin{problem}
如果函数 \(f(x)=\sin(\pi x+\theta)\)（\(0<\theta<2\pi\)）的最小正周期是 \(T\)，且当 \(x=2\) 时取得最大值，那么（\(\quad\)）

\choices
    {\(T=2\)，\(\theta=\dfrac{\pi}{2}\)}
    {\(T=1\)，\(\theta=\pi\)}
    {\(T=2\)，\(\theta=\pi\)}
    {\(T=1\)，\(\theta=\dfrac{\pi}{2}\)}
\end{problem}

\begin{answer}
A
\end{answer}

\begin{solution}
函数 \(\sin(\pi x+\theta)\) 的最小正周期为 \(T=\dfrac{2\pi}{\pi}=2\). 当 \(x=2\) 时函数取得最大值，因此 \(\sin(2\pi+\theta)=1\)，从而 \(\theta=\dfrac{\pi}{2}+2k\pi\). 结合 \(0<\theta<2\pi\)，得 \(\theta=\dfrac{\pi}{2}\)，故选 A.
\end{solution}

\begin{problem}
设 \(abc\ne0\)，“\(ac>0\)”是“曲线 \(ax^2+by^2=c\) 为椭圆”的（\(\quad\)）

\choices
    {充分非必要条件}
    {必要非充分条件}
    {充分必要条件}
    {既非充分又非必要条件}
\end{problem}

\begin{answer}
B
\end{answer}

\begin{solution}
若曲线 \(ax^2+by^2=c\) 是椭圆，则 \(a\)、\(b\)、\(c\) 必须同号，特别有 \(ac>0\). 反之，\(ac>0\) 不能保证是椭圆，例如取 \(a=c=1\)，\(b=-1\)，曲线为 \(x^2-y^2=1\)，它是双曲线而不是椭圆. 因此 \(ac>0\) 是必要非充分条件，故选 B.
\end{solution}

\begin{problem}
直线 \(x+\sqrt{3}y-2=0\) 被圆 \((x-1)^2+y^2=1\) 所截得的线段的长为（\(\quad\)）

\choices
    {\(1\)}
    {\(\sqrt{2}\)}
    {\(\sqrt{3}\)}
    {\(2\)}
\end{problem}

\begin{answer}
C
\end{answer}

\begin{solution}
圆心为 \(C(1,0)\)，半径为 \(1\). 圆心到直线的距离为
\[
d=\frac{|1+\sqrt{3}\cdot0-2|}{\sqrt{1+(\sqrt{3})^2}}=\frac12.
\]
因此直线截圆所得弦长为
\[
2\sqrt{1^2-d^2}=2\sqrt{1-\frac14}=\sqrt{3},
\]
故选 C.
\end{solution}

\begin{problem}
在 \(\triangle ABC\) 中，已知 \(2\sin A\cos B=\sin C\)，那么 \(\triangle ABC\) 一定是（\(\quad\)）

\choices
    {直角三角形}
    {等腰三角形}
    {等腰直角三角形}
    {正三角形}
\end{problem}

\begin{answer}
B
\end{answer}

\begin{solution}
由三角形内角和 \(C=\pi-A-B\)，得 \(\sin C=\sin(A+B)=\sin A\cos B+\cos A\sin B\). 代入已知等式，得到 \(\sin A\cos B=\cos A\sin B\)，即 \(\sin(A-B)=0\). 由于 \(A,B\in(0,\pi)\)，所以 \(A-B\in(-\pi,\pi)\)，只能有 \(A=B\). 因此三角形为等腰三角形，故选 B.
\end{solution}

\begin{problem}
若不等式 \((-1)^na<2+\dfrac{(-1)^{n+1}}{n}\) 对于任意正整数 \(n\) 恒成立，则实数 \(a\) 的取值范围是（\(\quad\)）

\choices
    {\(\left[-2,\dfrac{3}{2}\right)\)}
    {\(\left(-2,\dfrac{3}{2}\right)\)}
    {\(\left[-3,\dfrac{3}{2}\right)\)}
    {\(\left(-3,\dfrac{3}{2}\right)\)}
\end{problem}

\begin{answer}
A
\end{answer}

\begin{solution}
当 \(n=2m\) 为偶数时，不等式化为 \(a<2-\dfrac{1}{2m}\). 其中最严格的条件由 \(m=1\) 给出，即 \(a<\dfrac32\)；满足此条件时所有偶数情形都成立.

当 \(n=2m-1\) 为奇数时，不等式化为 \(a>-2-\dfrac{1}{2m-1}\). 若 \(a\geq-2\)，则对任意正整数 \(m\) 都成立；若 \(a<-2\)，由于 \(-2-\dfrac{1}{2m-1}\) 可以任意接近 \(-2\)，取充分大的 \(m\) 即可使该不等式不成立. 因此奇数情形等价于 \(a\geq-2\).

综上，\(a\) 的取值范围为 \(\left[-2,\dfrac32\right)\)，故选 A.
\end{solution}

\section{填空题}
\begin{problem}
\(\displaystyle\lim_{n\to\infty}\dfrac{n^2}{2n^2-3}=\)\fillinblank{}.
\end{problem}

\begin{answer}
\(\dfrac12\)
\end{answer}

\begin{solution}
分子、分母同除以 \(n^2\)，得 \(\displaystyle\frac{n^2}{2n^2-3}=\frac{1}{2-\frac{3}{n^2}}\). 当 \(n\to\infty\) 时，\(\frac{3}{n^2}\to0\)，所以原极限为 \(\dfrac12\).
\end{solution}

\begin{problem}
椭圆 \(\dfrac{x^2}{25}+\dfrac{y^2}{9}=1\) 的离心率是\fillinblank{}，准线方程是\fillinblank{}.
\end{problem}

\begin{answer}
\(\dfrac45\)，\(x=\pm\dfrac{25}{4}\)
\end{answer}

\begin{solution}
椭圆的长半轴为 \(a=5\)，短半轴为 \(b=3\)，所以焦半距满足 \(c^2=a^2-b^2=25-9=16\)，即 \(c=4\). 因此离心率为 \(e=\dfrac{c}{a}=\dfrac45\)，准线方程为 \(x=\pm\dfrac{a}{e}=\pm\dfrac{25}{4}\).
\end{solution}

\begin{problem}
已知 \(\sin\dfrac{\theta}{2}+\cos\dfrac{\theta}{2}=\dfrac{2\sqrt{3}}{3}\)，那么 \(\sin\theta\) 的值为\fillinblank{}，\(\cos2\theta\) 的值为\fillinblank{}.
\end{problem}

\begin{answer}
\(\dfrac13\)，\(\dfrac79\)
\end{answer}

\begin{solution}
将已知等式两边平方，得 \(\sin^2\dfrac\theta2+\cos^2\dfrac\theta2+2\sin\dfrac\theta2\cos\dfrac\theta2=\dfrac43\)，即 \(1+\sin\theta=\dfrac43\)，所以 \(\sin\theta=\dfrac13\). 再由二倍角公式，\(\cos2\theta=1-2\sin^2\theta=1-2\left(\dfrac13\right)^2=\dfrac79\).
\end{solution}

\begin{problem}
如图，正方体 \(ABCD-A_1B_1C_1D_1\) 的棱长为 \(a\)，将该正方体沿对角面 \(BB_1D_1D\) 切成两块，再将这两块拼接成一个不是正方体的四棱柱，那么所得四棱柱的全面积为\fillinblank{}.

\begin{center}
\bitmapfigure[max width=0.42\linewidth]{img/2005/beijing_liberal_spring/q12_fig1.png}
\end{center}
\end{problem}

\begin{answer}
\((4+2\sqrt{2})a^2\)
\end{answer}

\begin{solution}
切开后得到两个全等的三棱柱，每个三棱柱的三角形底面是直角等腰三角形，直角边均为 \(a\)，斜边为 \(a\sqrt2\). 一个三棱柱的两个三角形底面面积之和为 \(a^2\)，三个侧面面积之和为 \(a^2+a^2+a^2\sqrt2=(2+\sqrt2)a^2\)，所以其表面积为 \((3+\sqrt2)a^2\).

拼接成四棱柱时，两块沿一个面积为 \(a^2\) 的矩形面相接，该矩形面的两份面积变为内部面积. 因此所得四棱柱的全面积为 \(2(3+\sqrt2)a^2-2a^2=(4+2\sqrt2)a^2\).
\end{solution}

\begin{problem}
从 \(0\)，\(1\)，\(2\)，\(3\) 这四个数中选三个不同的数作为函数 \(f(x)=ax^2+bx+c\) 的系数，可组成不同的一次函数有\fillinblank{}个，不同的二次函数有\fillinblank{}个.（用数字作答）
\end{problem}

\begin{answer}
\(6\)，\(18\)
\end{answer}

\begin{solution}
一次函数要求 \(a=0\)，再从 \(1,2,3\) 中选两个不同的数分别作为 \(b,c\)，有 \(3\times2=6\) 个.

二次函数要求 \(a\ne0\). 先从 \(1,2,3\) 中选 \(a\)，有 \(3\) 种；再从剩下的三个数中有序选取 \(b,c\)，有 \(3\times2\) 种. 因此不同的二次函数有 \(3\times3\times2=18\) 个.
\end{solution}

\begin{problem}
若关于 \(x\) 的不等式 \(x^2-ax-a>0\) 的解集为 \((-\infty,+\infty)\)，则实数 \(a\) 的取值范围是\fillinblank{}.
\end{problem}

\begin{answer}
\((-4,0)\)
\end{answer}

\begin{solution}
二次函数 \(x^2-ax-a\) 的开口向上. 要使它对任意实数 \(x\) 都严格大于零，必须且只需其判别式小于零，即 \(\Delta=(-a)^2-4\cdot1\cdot(-a)=a^2+4a<0\). 解得 \(-4<a<0\)，故实数 \(a\) 的取值范围为 \((-4,0)\).
\end{solution}

\section{解答题}
\begin{problem}
设函数 \(f(x)=\lg(2x-3)\) 的定义域为集合 \(M\)，函数 \(g(x)=\sqrt{(x-3)(x-1)}\) 的定义域为集合 \(N\). 求：

\begin{enumerate}
\item 集合 \(M\)，\(N\)；
\item 集合 \(M\cap N\)，\(M\cup N\).
\end{enumerate}
\end{problem}

\begin{solution}
\begin{enumerate}
\item 因为对数的真数必须为正，所以 \(2x-3>0\)，得 \(M=\left(\dfrac32,+\infty\right)\). 因为根式内的式子必须非负，所以 \((x-3)(x-1)\geq0\)，得 \(N=(-\infty,1]\cup[3,+\infty)\).
\item 由上述两个集合直接得到 \(M\cap N=[3,+\infty)\)，\(M\cup N=(-\infty,1]\cup\left(\dfrac32,+\infty\right)\).
\end{enumerate}
\end{solution}

\begin{problem}
如图，正三棱锥 \(S-ABC\) 中，底面边长是 \(3\)，棱锥的侧面积等于底面积的 \(2\) 倍，\(M\) 是 \(BC\) 的中点. 求：

\begin{enumerate}
\item \(\dfrac{AM}{SM}\) 的值；
\item 二面角 \(S-BC-A\) 的大小；
\item 正三棱锥 \(S-ABC\) 的体积.
\end{enumerate}

\begin{center}
\bitmapfigure[max width=0.48\linewidth]{img/2005/beijing_liberal_spring/q16_fig1.png}
\end{center}
\end{problem}

\begin{solution}
设正三棱锥的顶点 \(S\) 在底面上的射影为正三角形中心 \(O\). 底面面积为 \(\dfrac{9\sqrt3}{4}\)，所以每个侧面的面积为 \(\dfrac13\cdot2\cdot\dfrac{9\sqrt3}{4}=\dfrac{3\sqrt3}{2}\).

\begin{enumerate}
\item 在等边三角形 \(ABC\) 中，\(AM=\dfrac{3\sqrt3}{2}\). 又因为 \(M\) 是侧面 \(SBC\) 的底边中点，\(SM\) 是该侧面的高，故 \(\dfrac12\cdot3\cdot SM=\dfrac{3\sqrt3}{2}\)，从而 \(SM=\sqrt3\). 因此 \(\dfrac{AM}{SM}=\dfrac32\).
\item 由正三角形的性质，\(A,O,M\) 共线，且 \(OM=\dfrac13AM=\dfrac{\sqrt3}{2}\)，\(AO=\dfrac23AM=\sqrt3\). 在直角三角形 \(SOM\) 中，\(SO=\sqrt{SM^2-OM^2}=\dfrac32\). 于是 \(AS^2=SO^2+AO^2=\dfrac{21}{4}\).

二面角 \(S-BC-A\) 的平面角是直线 \(SM\) 与 \(AM\) 的夹角. 在三角形 \(SAM\) 中，设该角为 \(\varphi\)，则
\[
\cos\varphi=\frac{SM^2+AM^2-AS^2}{2\cdot SM\cdot AM}=\frac{3+\frac{27}{4}-\frac{21}{4}}{2\cdot\sqrt3\cdot\frac{3\sqrt3}{2}}=\frac12.
\]
所以 \(\varphi=60^\circ\).
\item 正三棱锥的体积为
\[
V=\frac13\cdot\frac{9\sqrt3}{4}\cdot SO=\frac13\cdot\frac{9\sqrt3}{4}\cdot\frac32=\frac{9\sqrt3}{8}.
\]
\end{enumerate}
\end{solution}

\begin{problem}
已知 \(\{a_n\}\) 是等比数列，\(a_1=2\)，\(a_4=54\)；\(\{b_n\}\) 是等差数列，\(b_1=2\)，\(b_1+b_2+b_3+b_4=a_1+a_2+a_3\).

\begin{enumerate}
\item 求数列 \(\{a_n\}\) 的通项公式及前 \(n\) 项和 \(S_n\) 的公式；
\item 求数列 \(\{b_n\}\) 的通项公式；
\item 设 \(U_n=b_1+b_4+b_7+\cdots+b_{3n-2}\)，其中 \(n=1,2,\cdots\)，求 \(U_{10}\) 的值.
\end{enumerate}
\end{problem}

\begin{solution}
\begin{enumerate}
\item 设等比数列 \(\{a_n\}\) 的公比为 \(q\)，则 \(a_4=a_1q^3\)，所以 \(2q^3=54\)，得 \(q=3\). 因此 \(a_n=2\cdot3^{n-1}\)，前 \(n\) 项和为 \(S_n=\dfrac{2(3^n-1)}{3-1}=3^n-1\).
\item 由第（1）问得 \(a_1+a_2+a_3=2+6+18=26\). 设等差数列 \(\{b_n\}\) 的公差为 \(d\)，则 \(b_1+b_2+b_3+b_4=\dfrac{4(2+2+3d)}{2}=8+6d\). 由题意 \(8+6d=26\)，得 \(d=3\)，所以 \(b_n=2+3(n-1)=3n-1\).
\item \(b_1,b_4,b_7,\ldots,b_{3n-2}\) 是首项为 \(b_1=2\)、公差为 \(3d=9\) 的等差数列. 其第十项为 \(b_{28}=3\cdot28-1=83\)，故
\[
U_{10}=\frac{10(2+83)}{2}=425.
\]
\end{enumerate}
\end{solution}

\begin{problem}
如图，\(O\) 为坐标原点，过点 \(P(2,0)\) 且斜率为 \(k\) 的直线 \(l\) 交抛物线 \(y^2=2x\) 于 \(M(x_1,y_1)\)，\(N(x_2,y_2)\) 两点.

\begin{enumerate}
\item 写出直线 \(l\) 的截距式方程；
\item 求 \(x_1x_2\) 与 \(y_1y_2\) 的值；
\item 求证：\(OM\perp ON\).
\end{enumerate}

\begin{center}
\bitmapfigure[max width=0.42\linewidth]{img/2005/beijing_liberal_spring/q18_fig1.png}
\end{center}
\end{problem}

\begin{solution}
因为直线 \(l\) 过 \(P(2,0)\) 且斜率为 \(k\)，所以 \(k\ne0\)，并且其方程为 \(y=k(x-2)\).

\begin{enumerate}
\item 直线在 \(x\) 轴上的截距为 \(2\)，在 \(y\) 轴上的截距为 \(-2k\)，所以截距式方程为 \(\dfrac{x}{2}-\dfrac{y}{2k}=1\).
\item 由直线方程得 \(x=2+\dfrac{y}{k}\). 代入抛物线方程，得 \(y^2=2\left(2+\dfrac{y}{k}\right)\)，即 \(y^2-\dfrac{2}{k}y-4=0\). 因此 \(y_1,y_2\) 是此方程的两个根，由韦达定理得 \(y_1y_2=-4\). 又因 \(x_i=\dfrac{y_i^2}{2}\)，所以 \(x_1x_2=\dfrac{y_1^2y_2^2}{4}=4\).
\item 直线 \(OM\)、\(ON\) 的斜率分别为 \(\dfrac{y_1}{x_1}=\dfrac{2}{y_1}\)、\(\dfrac{y_2}{x_2}=\dfrac{2}{y_2}\)，它们的乘积为 \(\dfrac{4}{y_1y_2}=-1\). 所以 \(OM\perp ON\).
\end{enumerate}
\end{solution}

\begin{problem}
经过长期观测得到：在交通繁忙的时段内，某公路段汽车的车流量 \(y\)（\(\text{千辆/小时}\)）与汽车的平均速度 \(v\)（\(\text{千米/小时}\)）之间的函数关系为：
\(y=\dfrac{920v}{v^2+3v+1600}\)（\(v>0\)）.

\begin{enumerate}
\item 在该时段内，当汽车的平均速度 \(v\) 为多少时，车流量最大？最大车流量为多少？（精确到 \(0.1\text{ 千辆/小时}\)）
\item 若要求在该时段内车流量超过 \(10\text{ 千辆/小时}\)，则汽车的平均速度应在什么范围内？
\end{enumerate}
\end{problem}

\begin{solution}
\begin{enumerate}
\item 当 \(v>0\) 时，\(v^2+3v+1600>0\)，且由基本不等式 \(v+\dfrac{1600}{v}\geq2\sqrt{1600}=80\)，得
\[
v^2+3v+1600=v\left(v+3+\frac{1600}{v}\right)\geq83v.
\]
因此
\[
y=\frac{920v}{v^2+3v+1600}\leq\frac{920}{83},
\]
当且仅当 \(v=\dfrac{1600}{v}\)，即 \(v=40\) 时取等号. 故平均速度为 \(40\text{ 千米/小时}\) 时车流量最大，最大车流量为 \(\dfrac{920}{83}\approx11.1\text{ 千辆/小时}\).
\item 因为分母为正，车流量超过 \(10\text{ 千辆/小时}\) 等价于
\[
\frac{920v}{v^2+3v+1600}>10,
\]
即 \(920v>10v^2+30v+16000\)，整理得 \(v^2-89v+1600<0\).
又 \(v^2-89v+1600=(v-25)(v-64)\)，所以 \(25<v<64\).
\end{enumerate}
\end{solution}

\begin{problem}
现有一组互不相同且从小到大排列的数据：\(a_0\)，\(a_1\)，\(a_2\)，\(a_3\)，\(a_4\)，\(a_5\)，其中 \(a_0=0\). 为提取反映数据间差异程度的某种指标，今对其进行如下加工：记 \(T=a_0+a_1+\cdots+a_5\)，\(x_n=\dfrac{n}{5}\)，\(y_n=\dfrac{1}{T}(a_0+a_1+\cdots+a_n)\)，作函数 \(y=f(x)\)，使其图象为逐点依次连接点 \(P_n(x_n,y_n)\)（\(n=0,1,2,\cdots,5\)）的折线.

\begin{enumerate}
\item 求 \(f(0)\) 和 \(f(1)\) 的值；
\item 设 \(P_{n-1}P_n\) 的斜率为 \(k_n\)（\(n=1,2,3,4,5\)），判断 \(k_1\)，\(k_2\)，\(k_3\)，\(k_4\)，\(k_5\) 的大小关系；
\item 证明：\(f(x_n)<x_n\)（\(n=1,2,3,4\)）.
\end{enumerate}
\end{problem}

\begin{solution}
记 \(R_n=a_0+a_1+\cdots+a_n\)，则 \(y_n=\dfrac{R_n}{T}\).

\begin{enumerate}
\item 因为 \(a_0=0\)，所以 \(P_0=(0,0)\)，从而 \(f(0)=0\). 又因为 \(R_5=T\)，所以 \(P_5=(1,1)\)，从而 \(f(1)=1\).
\item 由相邻两点的坐标，得
\[
k_n=\frac{y_n-y_{n-1}}{x_n-x_{n-1}}=\frac{\frac{a_n}{T}}{\frac15}=\frac{5a_n}{T}.
\]
由于 \(a_1<a_2<a_3<a_4<a_5\)，所以 \(k_1<k_2<k_3<k_4<k_5\).
\item 对 \(n=1,2,3,4\)，由于 \(a_0=0\) 且 \(a_1<a_2<\cdots<a_5\)，有 \(R_n\leq na_n\)，并且 \(a_{n+1}+\cdots+a_5>(5-n)a_n\). 因此
\[
nT-5R_n=n(T-R_n)-(5-n)R_n>n(5-n)a_n-(5-n)R_n\geq0.
\]
于是 \(nT>5R_n\)，从而 \(f(x_n)=\dfrac{R_n}{T}<\dfrac{n}{5}=x_n\).
\end{enumerate}
\end{solution}
